\section{The fourfold: integral data, periods, and topology}
\label{sec:fourfold}

We use column vectors.  The matrices $A,B$ in the supplied data act on
cohomology; $M_1=A^t$, $M_2=B^t$ act on
$\Lambda=\mathbb Z^6=H_1(T^6;\mathbb Z)$, and
$M_0=(M_1M_2)^{-1}$.  Positive based meridians satisfy
$g_1g_2g_0=1$.  The exterior-square basis of $\bigwedge^2\mathbb Z^4$
is $(01,02,03,12,13,23)$.  All assertions about an integral matrix below
are checked over $\mathbb Z$ or $\mathbb Q$, without floating-point
arithmetic.  The files in \texttt{code/part1/} contain the complete
matrices, including the changes of basis in every Smith computation.

\begin{statement}{P1.1}{Proved}
The finite homology monodromies have orders $5,4$, respectively, and
characteristic polynomials
$(T-1)^2\Phi _5(T)$ and $(T-1)^2(T+1)^2(T^2+1)$.
For $N=M_0-I$ we have $N^2=0$, $\operatorname{rank}N=3$, and
nonzero Smith factors $1,1,1$.  Moreover
\[
 \Lambda^{\langle M_1,M_2\rangle}=0,\qquad
 \Lambda_{\langle M_1,M_2\rangle}=\mathbb Z,
\]
where the quotient homomorphism is $\ell(x)=x_0$.
On the cohomology lattice the corresponding groups are $\mathbb Ze_0$
and $\mathbb Z/2$.  Every nonidentity power of each specified affine
action $\tau_i$ is free.
\end{statement}
\begin{proof}
Exact multiplication gives
\[
N=\begin{pmatrix}
0&0&0&0&0&0\\0&0&0&1&0&-1\\0&0&0&0&1&1\\
-1&0&0&0&0&0\\1&0&0&0&0&0\\-1&0&0&0&0&0
\end{pmatrix},\qquad
S=\begin{pmatrix}
0&0&0&1&0&0\\0&1&0&0&0&0\\0&0&1&0&0&0\\
-1&0&0&0&1&0\\1&0&0&0&0&1\\-1&0&0&0&0&0
\end{pmatrix}.
\]
The matrix $S$ is unimodular and
$S^{-1}M_0S=\left(\begin{smallmatrix}I_3&I_3\\0&I_3\end{smallmatrix}\right)$.
The columns of $S$ are henceforth denoted
$(a_1,a_2,a_3,b_1,b_2,b_3)$, so that $Nb_j=a_j$.
This is also an integral certificate of the Smith assertion.
The Hermite basis of
$(M_1-I)\Lambda+(M_2-I)\Lambda$ is $(e_1,\ldots,e_5)$;
the stacked matrix of the two fixed-vector equations has rank six.
Transposing the matrices gives a fixed lattice $\mathbb Ze_0$ and
coinvariant Smith factors $(1,1,1,1,1,2)$.  The pairing
$\Lambda^\vee\times\Lambda\longrightarrow\mathbb Z$ identifies the
fixed cohomology vector $e_0$ with the covector $\ell$.

We have $M_iv_i=v_i$, $\ell(v_1)=1$, and $\ell(v_2)=-1$.
Thus $\tau_i^k(x)=M_i^kx+kv_i/m_i$, and $\tau_i^{m_i}=1$ on the
torus.  A fixed point of $\tau_i^k$ would give
$k\ell(v_i)/m_i=0$ in $\mathbb R/\mathbb Z$.  For
$0<k<m_i$ this is impossible.  The order and characteristic-polynomial
computations, and the full fixed-lattice calculations, are reproduced by
\texttt{tuple\_check.py}; its output records every power tested.
\end{proof}

\begin{statement}{P1.2}{Proved}
The lower blocks of $A,B$ are precisely the quotient representations of
$\bigwedge^2(C_f^{-1}),\bigwedge^2C_g$ in the supplied basis $Q$.
The integral quotient is the dual of the primitive exterior-square
lattice.  The natural map from the primitive lattice into that quotient
has cokernel $\mathbb Z/2$; these two integral representations are not
isomorphic.  Their rationalizations are isomorphic by that same map.
\end{statement}
\begin{proof}
The two matrices obtained from their two-by-two minors are
\[
Q^{-1}\bigwedge^2(C_f^{-1})Q=
\begin{pmatrix}
1&-1&0&1&0&0\\0&1&-1&-1&1&0\\0&1&0&-1&0&0\\
0&2&-1&-1&0&1\\0&2&0&-1&0&0\\0&1&1&-1&0&0
\end{pmatrix},
\]
\[
Q^{-1}\bigwedge^2C_gQ=
\begin{pmatrix}
1&0&1&0&0&0\\0&0&-1&0&1&0\\0&0&-1&0&0&1\\
0&0&-1&0&0&0\\0&1&-1&0&0&0\\0&0&-1&1&0&0
\end{pmatrix}.
\]
Their lower five-by-five blocks are the asserted blocks, entry by entry.
Put $L=\bigwedge^2\mathbb Z^4$, let $(\ ,\ )$ be its unimodular wedge
pairing, and set $\omega=(1,1,1,1,1,1)^t$.
Then $(\omega,\omega)=2$ and the primitive sublattice
$P_{\mathbb Z}=\omega^\perp$ is the kernel of
$(1,-1,1,1,-1,1)$.  In the basis with columns
\[
\begin{pmatrix}
1&-1&-1&1&-1\\1&0&0&0&0\\0&1&0&0&0\\
0&0&1&0&0\\0&0&0&1&0\\0&0&0&0&1
\end{pmatrix},
\]
the inclusion followed by projection
$j:P_{\mathbb Z}\longrightarrow L/\mathbb Z\omega$ has matrix
\[
j=\begin{pmatrix}
0&1&1&-1&1\\-1&2&1&-1&1\\-1&1&2&-1&1\\
-1&1&1&0&1\\-1&1&1&-1&2
\end{pmatrix},\qquad \det j=2.
\]
Its Smith factors are $(1,1,1,1,2)$.  The unimodular pairing on $L$
identifies $L/\mathbb Z\omega$ with $P_{\mathbb Z}^\vee$;
under this identification $j$ is the pairing map.

To rule out a different integral intertwiner, we solve the two
simultaneous bilinear-form invariance equations.  Each solution space
over $\mathbb Q$ has dimension one.  The primitive integral forms have
determinants $-2$ on $P_{\mathbb Z}$ and $-16$ on the quotient.
A unimodular intertwiner would pull one primitive form to a rational
multiple of the other.  Having coefficient content one on both sides
forces that multiple to be $1$ or $-1$, which contradicts the different
determinants.  All these linear equations and their solutions occur in
\texttt{primitive\_lattice.py/.out}.  Thus the rational wording in the
data is correct, whereas the stronger integral identification requested
in item 1.4 has precisely the defect just exhibited.
\end{proof}

\begin{statement}{P1.3}{Proved}
The two finite filling pieces have first integral homology
$\mathbb Z^2$ and $\mathbb Z^2\oplus\mathbb Z/2$.
\end{statement}
\begin{proof}
For $K_i=T^6/\langle\tau_i\rangle$, the universal covering deck group
has commuting translations $x^u$, $u\in\Lambda$, and a generator $a_i$,
with the complete relations
\[
a_ix^ua_i^{-1}=x^{M_iu},\qquad a_i^{m_i}=x^{v_i}.
\]
Indeed the affine lifts satisfy these relations, and they reduce every
word uniquely to $x^ua_i^r$, $0\leq r<m_i$.
The abelianization relation matrix has columns
$((M_i-I)e_j,0)$ and $(-v_i,m_i)$.
Smith reduction of these explicit seven-by-seven matrices gives the
groups above.  The disk quotient
$(T^6\times\Delta)/\langle(\tau_i,\zeta_{m_i})\rangle$
retracts onto $K_i$ by equivariant radial contraction.  Freeness supplies
ordinary smooth quotient charts by disjoint translates of small balls.
\end{proof}

\begin{statement}{P1.4}{Proved}
The genus-two pencil has one nodal fibre away from $0,\infty$, at
$\lambda_*=-256/3125$.  Good reduction at $0,\infty$ requires root
orders $8,5$.  On this root orbifold its Hodge bundle is
$\mathcal H=T^3\oplus T$, where
\[
T=\mathcal O\bigl((2/5)[\infty]-(3/8)[0]\bigr),\qquad
\deg T=1/40,
\]
and $T^3$ is the line generated meromorphically by $dx/y$.
\end{statement}
\begin{proof}
The discriminant of $x^5+\lambda^3x+\lambda^4$ is
$\lambda^{15}(3125\lambda+256)$.
At $\lambda_*$ the double root is $64/625$;
the second $x$-derivative and first $\lambda$-derivative there are
nonzero.  The remaining roots are simple.  The odd-degree
hyperelliptic compactification is smooth at its branch point at
infinity.  This gives one node, a smooth total surface at that node,
and genus-one normalization.

At zero set $\lambda=u^8$, $x=u^6X$, $y=u^{15}Y$; at infinity set
$\lambda=u^{-5}$, $x=u^{-4}X$, $y=u^{-10}Y$.
The resulting smooth equations and differentials are
\[
\begin{array}{c|c|c|c}
 &\text{equation}&dx/y&x\,dx/y\\\hline
0&Y^2=X^5+X+u^2&u^{-9}dX/Y&u^{-3}X\,dX/Y\\
\infty&Y^2=X^5+uX+1&u^6dX/Y&u^2X\,dX/Y.
\end{array}
\]
Their central polynomials are square-free.  The deck maps at zero and
infinity are $(X,Y)\mapsto(iX,\zeta_8Y)$ and
$(X,Y)\mapsto(\zeta_5^{-1}X,Y)$, respectively.
The eigenvalues on $H^{1,0}$ are $(\zeta_8,\zeta_8^3)$ and
$(\zeta_5^{-1},\zeta_5^{-2})$.  Including conjugates gives characteristic
polynomials $T^4+1$ and $\Phi_5$, proving minimality of the root orders.
The fourth power at zero is the hyperelliptic involution, acting as
$-I$ on $H^1$ and as $I$ on its exterior square.

The two meromorphic differential generators have orbifold divisors
$(6/5)[\infty]-(9/8)[0]$ and
$(2/5)[\infty]-(3/8)[0]$.
At the node they are independent sections of the dualizing Hodge
extension: one has opposite nonzero residues at the two points of the
normalization, and the linear combination vanishing at the node is a
regular nonzero elliptic differential.  Thus there is no additional
divisor at $\lambda_*$.  These computations prove the asserted direct
sum and its degrees.  The discriminant and all substitutions are
checked in \texttt{pencil\_charts.py/.out}.
\end{proof}

\begin{statement}{P1.5}{Refuted}
The formula in the prompt with $P$ the image of
$\ell_0\wedge\mathcal H$, and $\mathcal H$ the rank-two Hodge bundle,
cannot produce a rank-three subbundle.  The required rank-three image
is instead the image of $\ell_0\wedge\mathcal H_{\mathrm{dR}}$, where
$\mathcal H_{\mathrm{dR}}$ is the full rank-four de Rham bundle.
\end{statement}
\begin{proof}
The kernel of $\ell_0\otimes\mathcal H\to\bigwedge^2\mathcal H$
is $\ell_0\otimes\ell_0$, so its image has rank one.  Projection and
addition of a map into the invariant line cannot increase its domain
rank.  For the corrected domain there is an injective holomorphic map
\[
\ell_0\otimes(\mathcal H_{\mathrm{dR}}/\ell_0)
 \longrightarrow\bigwedge^2\mathcal H_{\mathrm{dR}},\quad
u\otimes[v]\longmapsto u\wedge v.
\]
Every nonzero image vector is decomposable and has wedge square zero;
$\omega^2\ne0$.  Projection modulo $\mathcal O\omega$ is therefore
injective on its rank-three image.  This specifies the comparison map
and the corrected bundle used below.
\end{proof}

\begin{statement}{P1.6}{Known}
There is an orientation-preserving diffeomorphism
\[
\bigl(3\mathbb HP^2\mathbin{\#}(-\mathbb HP^2)\bigr)\mathbin{\#}\Sigma^8
 \cong 3\mathbb HP^2\mathbin{\#}(-\mathbb HP^2)
\]
for every homotopy eight-sphere $\Sigma^8$.
\end{statement}
\begin{proof}
Kramer--Stolz, \emph{A diffeomorphism classification of manifolds which
are like projective planes}, J. Differential Geometry 77 (2007),
Theorem A, applies to a smooth simply connected $2m$-manifold,
$2m\ne4$, with integral homology $\mathbb Z\oplus\mathbb Z\oplus\mathbb Z$.
Its conclusion is absorption of every homotopy $2m$-sphere by connected
sum.  The cell structure of $\mathbb HP^2$ in dimensions $0,4,8$
verifies all hypotheses.  The resulting diffeomorphism
$\mathbb HP^2\#\Sigma^8\to\mathbb HP^2$ preserves orientation, since
both signatures are $+1$, whereas reversing orientation negates
signature.  Taking its oriented connected sum with the identity on
$2\mathbb HP^2\#(-\mathbb HP^2)$ proves the displayed diffeomorphism.
We read the theorem in the publisher PDF at
\url{https://archive.ymsc.tsinghua.edu.cn/pacm_download/21/694-JDG-77-2-A1-kramer.pdf}.
\end{proof}

\begin{statement}{P1.7}{Proved}
The integral $H^1$ monodromy of the pencil has a companion marking with
finite pair $(C_f^{-1},C_g)$, and cusp product $C_fC_g^{-1}$.
Its exterior-square quotient is consequently the integral quotient
representation appearing in P1.2.
\end{statement}
\begin{proof}
Put $z=\lambda/\lambda_*$ and use the based-meridian convention
$g_\infty g_1g_0=1$ at $z=1/2$ of Beukers--Heckman,
Definition 3.1.  The two finite spectra are disjoint.  The nodal
monodromy is a rank-one Picard--Lefschetz transformation with primitive
vanishing cycle $\delta$ and coefficient $\pm1$.
Here primitivity follows by taking a transverse integral curve to the
nonseparating vanishing circle; coefficient one follows by turning the
annulus $xy=z-1$ once positively.  The sign depends only on the order
of the pairing and on homology versus cohomology.

Beukers--Heckman, Invent. Math. 95 (1989), Theorem 3.5 (Levelt), gives
the companion triple $(C_f,C_f^{-1}C_g,C_g^{-1})$.
The proof is defined over $\mathbb Q$ in this case: the intersection of
the three appropriate translates of the rational rank-three kernel
contains a rational vector whose four successive translates form a
basis.  The disjoint spectra imply irreducibility, since a common
proper invariant subspace or its quotient on which the rank-one
transformation is trivial would give a shared eigenvalue.  This
verifies the field of the resulting intertwiner as well as the
hypotheses of the theorem.

The integral upgrade uses the exact matrices
\[
d=(1,1,1,0)^t,\qquad
J=\begin{pmatrix}
0&-1&1&-1\\1&0&-1&1\\-1&1&0&-1\\1&-1&1&0
\end{pmatrix}.
\]
They satisfy $C_f^tJC_f=C_g^tJC_g=J$, $\det J=1$, and
$C_f^{-1}C_g=I-d(d^tJ)$.
The orbit $(d,C_fd,C_f^2d,C_f^3d)$ has alternating pairing
\[
\begin{pmatrix}0&-1&0&0\\1&0&-1&0\\0&1&0&-1\\0&0&1&0\end{pmatrix},
\]
of determinant one.  Under a rational intertwiner let $d$ map to
$a\delta$, and let the pulled-back intersection form be $cJ$.
The space of simultaneous invariant alternating forms has dimension
one by the exact invariance equations.  Comparing the rank-one
transvection formulas yields $c/a^2=\pm1$.
It follows that the four integral cycles
$\delta,H_\infty\delta,H_\infty^2\delta,H_\infty^3\delta$ have the
displayed unimodular pairing, with an overall sign.  Their sublattice
has index one in the unimodular integral $H^1$, since its determinant
is the square of that index.  Rescaling the intertwiner by $a$ thus
gives an integral marking.

The antidiagonal reversal matrix $R$ satisfies
$RC_fR^{-1}=C_f^{-1}$ and $RC_g^{-1}R^{-1}=C_g$.
Consequently the cusp product is $(C_gC_f^{-1})^{-1}$, matching
$(BA)^{-1}$ on the supplied cohomology system.  Its transpose is the
specified $M_0$.  The matrices and pairing calculation are recorded in
\texttt{companion\_markings.py/.out}.  The verified primary source is
\url{https://math.uchicago.edu/~drinfeld/Hypergeometric_motives/Beukers-Heckman.pdf},
pp.~331--332.
\end{proof}

\begin{statement}{P1.8}{Proved}
With the corrected rank-three bundle $P$ of P1.5, the supplied rank-six
system admits a holomorphic lifted family
$F_c=(\sigma+ic\alpha e)(P)$ on the punctured base.
For a specified splitting $\sigma$ there is a real analytic function
$h$ with finite attained minimum $h_{\min}$, and the exact domain of
real parameters giving $F_c\cap\overline{F_c}=0$ everywhere is
$c>-h_{\min}$.
At the cusp its periods have the form
$Z(q)=(\log q/(2\pi i))I_3+C(q)$ with $C$ holomorphic and bounded.
The construction extends to the two finite root charts and descends to
orders $5,4$ on the rank-six system.
\end{statement}
\begin{proof}
Work first on the root orbifold with orders $5,8$, leaving the nodal
point as a logarithmic point.  The symplectic pairing and the Hodge
subbundle give an exact sequence
$0\to\mathcal H\to\mathcal H_{\mathrm{dR}}\to\mathcal H^\vee\to0$.
This remains exact at the node by the residue pairing in its plumbing
chart.  With $\ell_0=T^3$ the corrected bundle fits into
\[
0\longrightarrow T^4\longrightarrow P
 \longrightarrow\mathcal O\oplus T^2\longrightarrow0.
\]
Contraction in $\ell_0\wedge\mathcal H_{\mathrm{dR}}$ defines
$\alpha:P\to\mathcal O$, using the unique lift supplied by P1.5;
it is surjective.  Thus it is not an asserted contraction on the whole
quotient by $\omega$, where that formula would not be well defined.

For $a=2,4,-2,-4$ the invariant coarse direct image of $T^a$ has degree
$-1$, by taking the floors of the two coefficients of its rational
divisor.  Finite-stabilizer invariants are exact in characteristic zero
by averaging.  Consequently $H^0(T^a)=H^1(T^a)=0$ for these exponents.
The two-chart Laurent-series computation for $\mathcal O(-1)$ proves
this vanishing directly.  The extension therefore splits, and
\[
P\cong T^4\oplus T^2\oplus\mathcal O,\qquad
H^1(P^\vee)=0,\quad H^0(P^\vee)=\mathbb C\alpha,
\quad H^0(P)=\mathbb Cg,
\]
where $\alpha(g)=1$.

Let $V$ be the logarithmic extension of the supplied rank-six flat
bundle.  At the cusp use its nilpotent exponential frame, and on each
finite root chart use the trivial flat extension with the deck action.
These charts give the exact bundle sequence
$0\to\mathcal Oe\to V\to W\to0$.
Its pullback to $P$ splits because $H^1(P^\vee)=0$.
All holomorphic splittings are $\sigma+z\alpha e$, $z\in\mathbb C$.
To remove this scalar choice we require the $e$-coordinate of
$\sigma(g)$ to be zero in the fixed flat frame at $z=1/2$.

Choose a local generator $u$ of $\ell_0$ and put
$r=i\pi(u\wedge\bar u)\in W_{\mathbb R}$, $A=\alpha(r)>0$;
the sign is that of the polarization.  The symplectic orthogonal
decomposition into $\langle u,\bar u\rangle$ and its two-dimensional
complement proves $P\cap\bar P=\mathbb Cr$.
Indeed $u\wedge v-\bar u\wedge w=t\omega$ has zero component in the
exterior square of that complement on the left and a nonzero such
component on the right unless $t=0$.  With $t=0$ its intersection is
the line generated by $u\wedge\bar u$.
Define $h$ by the exact equality
\[
\sigma(r)-\overline{\sigma(r)}=2iAh e.
\]
Both $r$ and $A$ scale by $|a|^2$ when $u$ scales by $a$, so this is
well defined.  Comparing the unique lifts of a multiple of $r$ proves
\[
F_c\cap\overline{F_c}=0\quad\Longleftrightarrow\quad h+c\ne0.
\]

The sign and growth at the cusp follow from the actual extension
matrices.  Put $D=C_fC_g^{-1}-I$ and $N_V=(BA)^{-1}-I$.
Only the zeroth column of $D$ is nonzero, and it is $(0,1,1,1)^t$;
the top row of $N_V$ is $(0,0,0,-1,1,-1)$.
For arbitrary $a,b\in\mathbb C^4$, two-by-two minor expansion gives
\[
(N_V)_{0,*}(0,\pi(a\wedge Db))
 =\operatorname{contr}(a\wedge Db)=a_0b_0,
 \qquad N_W\pi(a\wedge Db)=0.
\]
These polynomial identities, with arbitrary symbols for the entries of
$a,b$, are checked in \texttt{cusp\_lift\_asymptotic.py}.
Write $u=\exp(sD)a(q)$, $s=\log q/(2\pi i)$ and $y=\operatorname{Im}s$.
The vector $a$ is holomorphic, and $a_0(0)\ne0$ because $dx/y$ has
nonzero nodal residue.  In the logarithmic frame we obtain
\[
\exp(-sN_W)r=i\pi(a\wedge\bar a)+2y\pi(a\wedge D\bar a),
\quad
A=i\operatorname{contr}(a\wedge\bar a)+2y|a_0|^2.
\]
The splitting is bounded in that frame.  On returning by
$\exp(sN_V)$, the only quadratic contribution to the imaginary
$e$-coordinate is $2|a_0|^2y^2$.  All other terms are $O(y)$,
uniformly for $\operatorname{Re}s$ in a fixed unit interval.
Dividing by $A=2|a_0|^2y+O(1)$ proves
\[
h(q)=y+O(1)=-\frac{\log|q|}{2\pi}+O(1).
\]
The function extends continuously over both finite root charts and
tends to $+\infty$ at the cusp.  Compactness away from that cusp gives
an attained minimum.  Connectedness and the intermediate value
theorem give $h(U)=[h_{\min},\infty)$, with finite-root points included
when testing the completed family.  Thus the domain is precisely the
displayed half-line.  Replacing $\sigma$ by $\sigma+ib\alpha e$ replaces
$h$ by $h+b$ and sends $F_c$ to $F_{c+b}$, proving the stated
comparison between choices.

At $q=0$ the map $N_V:F_{\log}\to\ker N_V$ is an isomorphism.
On the quotient its rank is two and its kernel is generated by
$\pi(a\wedge Da)$; the displayed scalar identity sends a lift of this
generator to $a_0^2e\ne0$.  In the integral Jordan marking, the
logarithmic three-plane is therefore a graph of a holomorphic matrix
$C(q)$; applying $\exp(sN_V)$ adds $sI_3$.  Evaluation on the homology
lattice gives the period matrix $(I_3,Z(q))$.
The fourth power of the order-eight deck transformation acts trivially
on $V$ and $P$, so the graph descends to the order-four chart.
The other chart has order five.  The affine translations are
holomorphic translations on the resulting tori, and their linear
parts preserve this graph by construction.
\end{proof}

\paragraph{Scope of P1.8.}
The definition of $h$ is exact and specifies its normalization, its
range, and its cusp growth.  We have not evaluated $h_{\min}$ or
supplied the two independent thirty-digit period calculations requested
in item 1.4.  These numerical requirements remain separately recorded;
they are not consequences of the existence proof.

\begin{statement}{P1.9}{Proved}
The periodic Kuhn fan, in the marking $S$, gives a smooth proper
holomorphic filling of the periods $Z(q)=sI_3+C(q)$ for a sufficiently
small disk.  Its central fibre is reduced with normal crossings.
Modulo the index-one period lattice it has respectively $1,7,12,6$
strata of dimensions $3,2,1,0$, each isomorphic to the corresponding
complex torus $(\mathbb C^*)^d$, and Euler number $6$.
An index-$d$ period sublattice gives $d$ times these stratum counts and
Euler number $6d$.
\end{statement}
\begin{proof}
Let $\Sigma$ consist of the cones over all simplices of the periodic
Kuhn triangulation at height one, together with their faces.  For an
ordered maximal simplex $v_0,\ldots,v_3$, the columns $(v_j,1)$ form a
unimodular four-by-four matrix.  If $u_j$ are its dual integral basis,
the corresponding chart has coordinates $x_j=\chi^{u_j}$ and is
$\mathbb C^4$.  The base character is the sum of the four $u_j$, so
\[
q=x_0x_1x_2x_3.
\]
On a common face we invert the coordinates of omitted rays and change
dual bases; these are mutually inverse integral monomial maps.  These
explicit charts and overlaps construct a smooth toric space over the
disk with a reduced normal-crossings zero fibre.

The period group acts by
\[
(z,q)\longmapsto
\bigl(q^\lambda\exp(2\pi iC(q)\lambda)z,q\bigr),
\qquad \lambda\in\mathbb Z^3.
\]
The monomial part takes the cone over $\sigma$ to the cone over
$\sigma+\lambda$; the unit part is torus multiplication holomorphic in
$q$, which extends over each chart.  These maps compose according to
addition in $\mathbb Z^3$.  On the zero fibre a fixed point would
require a nonzero translation to stabilize a bounded simplex, which
is impossible.  For small nonzero $q$, the imaginary part of the
period matrix is $\operatorname{Im}(s)I_3+\operatorname{Im}C(q)$ and
is invertible; hence the action is free there too.

We give a compactness check for the infinite-chart construction.  Put
$u=\log|z|/\log|q|$.  A period changes $u$ by
$(I+E(q))\lambda$, where $E(q)=O(1/|\log|q||)$ uniformly.
On a small disk $(I+E)^{-1}$ has norm at most two.  Reduction of
$(I+E)^{-1}u$ modulo $\mathbb Z^3$ therefore gives a representative
with $u$ in a fixed bounded region.  Only finitely many simplices
meet that region.  In a simplex containing $u$, the numbers
$\log|x_j|/\log|q|$ are its barycentric coordinates; thus
$|x_j|\leq1$.  Finitely many closed polydisks in these charts give
compact representatives above a smaller closed base disk.  Every
central point is included in such representatives by translating its
finite star into the same bounded region and then using its complete
toric component.

The action is proper.  For compact coordinate neighborhoods approaching
the central fibre, only finitely many cones and their faces occur;
their translates can meet only when the translation is bounded.
Equivalently, the valuation estimate just given and the uniform bound
on $(I+E)^{-1}$ bound the translations for a sequence with $q\to0$.
On a closed annulus away from zero, properness is the usual
properness of a full lattice action.  These two checks give finitely
many intersecting translates of any pair of compact neighborhoods.
The quotient is therefore Hausdorff, is smooth by freeness, and is
proper over the disk by the compact representatives.  The map is flat
because $q$ is a nonzero divisor in each smooth chart.  Exponentiation
$z_j=\exp(2\pi i\zeta_j)$ identifies the nonzero fibres with
$\mathbb C^3/(\mathbb Z^3+Z(q)\mathbb Z^3)$.

A $p$-simplex supplies a stratum $(\mathbb C^*)^{3-p}$.
Enumeration of all faces of the six Kuhn tetrahedra, retaining their
translation identifications, gives $(1,7,12,6)$.
Additivity of compactly supported Euler characteristic and
$\chi_c(\mathbb C^*)=0$ leave six point contributions.
Inclusion of an index-$d$ period sublattice induces a finite covering
of these toric quotients.  Each simplex orbit has exactly $d$ lifts,
which proves the stated index dependence by an actual comparison map.
\end{proof}

\begin{statement}{P1.10}{Proved}
For every real parameter $c>-h_{\min}$ the corrected construction gives
a compact smooth complex fourfold $M_c$ and a proper holomorphic map
to $\mathbb P^1$ with the two specified multiple fibres and the
index-one Kuhn cusp.  With its complex orientation,
$\pi_1(M_c)=1$ and $\chi(M_c)=6$.
\end{statement}
\begin{proof}
Evaluation of $\Lambda$ on $F_c$ maps $\Lambda_{\mathbb R}$
isomorphically onto the underlying real bundle of $F_c^\vee$:
the kernel is zero exactly by $F_c\cap\bar F_c=0$.
The image is consequently a full discrete lattice.  Quotient by its
holomorphic period sections gives a holomorphic torus bundle on the
nonspecial base.

On a finite root chart $u^{m_i}=t_i$, divide the pullback torus bundle by
\[
(x,u)\longmapsto(M_ix+v_i/m_i,\zeta_{m_i}u).
\]
Freeness proves smoothness, including the central fibre.  The special
divisor is $m_i$ times the reduced free quotient.  On the punctured
root disk translation by
$-(\log u/(2\pi i))v_i$, interpreted by its period section, identifies
this affine quotient with the original Jacobian bundle.  A change of
logarithm gives an integral lattice period; replacing $u$ by
$\zeta_{m_i}u$ cancels exactly the affine term $v_i/m_i$.
This proves the holomorphic gluing and its sign.
At the cusp use P1.9 and the exponential period map.  The zero section
$\zeta=0$ extends to $z=1$ in the dense torus of the component marked
by the vertex zero.  The resulting map is proper by a finite compact
cover of the base by the three local pieces and a compact remainder
with torus fibres.  All total-space charts and their overlaps have
already been specified, so this produces a smooth compact complex
fourfold, and in particular an almost complex structure.

For completeness we record every van Kampen relation.  The main
bundle has its zero section and fundamental group
$\Lambda\rtimes\langle g_1,g_2\rangle$.
The finite pieces impose
$g_i^{m_i}=v_i$ and $g_ix^ug_i^{-1}=x^{M_iu}$; the logarithmic gluing
above supplies these signs.  The cusp zero section kills
$g_0=(g_1g_2)^{-1}$ and the cusp collapse kills $\operatorname{im}N$.
The latter assertion follows directly from the toric moment model in
P1.11: angular cycles $a_j$ collapse and the base cycles $b_j$ survive.
Its fundamental group is $\mathbb Z^3$, since the simply connected
complete toric components and their connected toric intersections have
dual complex $\mathbb R^3/\mathbb Z^3$; alternatively, van Kampen on
the moment cells kills the angular loops generated by the primitive
facet normals and retains exactly the three base loops.
Thus the full presentation is
\[
\left\langle \Lambda,g_1,g_2\ \middle|\
 [\Lambda,\Lambda]=1,\quad
 g_ix^ug_i^{-1}=x^{M_iu},\quad
 g_1^5=v_1,\ g_2^4=v_2,\ g_1g_2=1,\ \operatorname{im}N=0
\right\rangle.
\]
The last relation is redundant after the preceding ones but records
the actual cusp map.  Since $g_2=g_1^{-1}$, the two induced actions are
inverse; their orders divide both five and four and are therefore one.
The remaining fibre group is $\mathbb Z$ through $\ell$.  With its
generator $z$, the additive relations are
$5g_1-z=0$, $4g_2+z=0$, $g_1+g_2=0$.
Their matrix
\[
\begin{pmatrix}5&0&-1\\0&4&1\\1&1&0\end{pmatrix}
\]
has determinant $-1$, proving triviality.
Euler additivity gives $\chi(M_c)=6$: smooth tori and finite free
torus quotients have Euler number zero, and P1.9 gives the remaining
contribution.
\end{proof}

\begin{statement}{P1.11}{Proved}
The index-one central cusp fibre $C$ has torsion-free integral homology
with ranks $(1,3,9,9,10,3,1)$.
The collapse $\kappa:T^6\to C$ has primitive image in every degree;
it is onto except in degree four, where its cokernel is $\mathbb Z$.
Its image ranks are $(1,3,9,9,9,3,1)$.
In every degree its kernel is the saturation of
$(\bigwedge^kM_0-I)\bigwedge^k\Lambda$.
\end{statement}
\begin{proof}
We first remove the unit matrix $C(q)$ topologically, by an explicit
deformation, without changing the integral marking.  Replace it by
$\theta C(q)$, $0\leq\theta\leq1$.  The action and properness estimates
of P1.9 are uniform in $\theta$.  Local quotient charts form a smooth
family of pairs, with the local equation $q=x_0x_1x_2x_3$ independent
of $\theta$.  Local lifts of $\partial/\partial\theta$, tangent to every
coordinate stratum and annihilating $q$, can be patched by a partition
of unity.  Over a smaller compact disk their flow exists for the whole
interval.  This gives a stratified diffeomorphism of the central
fibres and a diffeomorphism of the nearby fibres, preserving the
period marking.  The same construction on the real oriented blowup
preserves the collapse map.  We may therefore calculate at $C(q)=0$.

Here is a concrete topological model.  Put $G=4I_3-\mathbf1\mathbf1^t$.
Its positive eigenvalues are $1,4,4$.  The Voronoi polytope in dual
coordinates is cut out by
\[
|v^ty|\leq \tfrac12v^tGv
 \quad (0\ne v\in\{0,1\}^3).
\]
The value on a binary vector with $k$ nonzero entries is $k(4-k)$.
These fourteen primitive facet normals give the star of a vertex in
the Kuhn triangulation.  One checks its Delaunay tetrahedra by the
successive chains
$0,e_i,e_i+e_j,e_1+e_2+e_3$ for the six permutations $(i,j,k)$;
solving their three perpendicular-bisector equations and the fourteen
inequalities verifies precisely their common faces.  The translates
of this Voronoi polytope by $G\mathbb Z^3$ tile $\mathbb R^3$.

For a simplex $\sigma$ let $D_\sigma$ be its dual cell in this tiling,
and let
$L_\sigma=\operatorname{span}_{\mathbb Z}\{v-v_0:v\in\sigma\}$.
The compact toric component with this normal fan is the quotient of
its moment polytope times the angular torus by the subtori generated
by the facet normals that meet a given face.  This description follows
on each affine chart by writing its coordinates in polar form;
at a zero coordinate its circle argument is forgotten.  Gluing the
faces and then taking the period quotient gives a continuous
homeomorphism from the compact quotient
\[
\mathcal Q=
\bigl((\mathbb R^3/G\mathbb Z^3)\times
       (\mathbb R^3/\mathbb Z^3)\bigr)/\!\sim\ \longrightarrow C,
\]
where over the relative interior of $D_\sigma$ the equivalence
relation collapses the subtorus $L_\sigma\otimes\mathbb R/\mathbb Z$.
The map is bijective by the polar-coordinate description on each
toric stratum, so compactness and Hausdorffness make it a homeomorphism.
Its angular identity defines a continuous section of the projection to
the base three-torus.

The actual nearby-fibre collapse is the quotient map in this model.
To see the comparison, perform the real oriented blowup of the
normal-crossings divisor in each chart.  Retain the arguments of all
$x_j$, with their sum fixed by $\arg q$, and then forget the arguments
of coordinates whose radii are zero.  At a face the forgotten
subtorus is exactly $L_\sigma$.  The uncollapsed arguments and the
dual moment coordinates give the product of the angular and base
three-tori above.  In a chart the passage to a small smooth fibre is
obtained by increasing the nonnegative radii continuously so their
product is $|q|$, keeping the arguments and $\arg q$ fixed; this is
invertible after fixing the transverse radial parameter.  On common
faces the integral monomial transformations give exactly the same
angular quotient.  Patching the radial parameters by positive
partitions of unity gives the global identification, while forgetting
the arguments is unchanged.  Thus this comparison fixes the map on
the marked angular cycles $a_j$ and base cycles $b_j$, not merely
the Betti numbers of a putative central fibre.

Filter $\mathcal Q$ by the dimension $d$ of its dual base cells.
Over the interior of the dual cell of a $p$-simplex, $d=3-p$ and the
angular quotient has lattice
$Q_\sigma=\mathbb Z^3/L_\sigma$ of rank $d$.
Excision for the cell and its boundary gives the relative chain
groups of the resulting homology spectral sequence:
\[
E^1_{d,q}=\bigoplus_{\dim\sigma=3-d}\bigwedge^qQ_\sigma.
\]
Orient dual cells by the orientation of the base and the ordered
simplex.  The differential is the signed quotient map for each
codimension-one coface.  More explicitly, with integral projection
matrices $Q_\sigma$ and chosen right inverses $S_\sigma$,
\[
D_{p,q}[\tau,\sigma]=
 \sum_{\sigma=\partial_j\tau\ {\rm mod}\ \mathbb Z^3}
 (-1)^j\bigwedge^q(Q_\tau S_\sigma).
\tag{1.11.1}
\]
All $L_\sigma$ are primitive because the tetrahedra are unimodular.
Thus these projections and right inverses are integral.  Formula
(1.11.1) specifies every integer entry, including multiple incidences
under translation identifications.  The program
\texttt{cusp\_complex.py} enumerates the $26$ simplex classes, all
$74$ codimension-one incidences, and exports every matrix, basis, and
Smith certificate to \texttt{cusp\_complex.json}.

The nonzero $E^2_{d,q}$ ranks are
\[
\begin{array}{c|rrrr}
q\backslash d&0&1&2&3\\\hline
0&1&3&3&1\\
1&0&6&8&3\\
2&0&0&7&3\\
3&0&0&0&1
\end{array}
\]
and every group is free.  Support leaves only the possible higher
differential $d_2:E^2_{3,0}\to E^2_{1,1}$.
The angular identity section splits the base projection as a map of
filtered spaces, so the base fundamental class survives.  This kills
that differential.  All subsequent differentials have zero source or
target.  The filtration extensions are extensions of free abelian
groups and therefore split, proving the asserted integral homology.

The product-to-quotient map induces on $E^1$ the exterior powers of
the quotient projections.  The program \texttt{cusp\_collapse.py}
computes the induced maps on the above Smith bases; all cokernels are
free, zero except for one generator in total degree four.
Surjectivity and primitivity transfer from the associated graded
groups through the finite filtration: successively lift a class from
each filtration quotient, and note that a filtration with free
quotients has no torsion.  This proves the corresponding assertions
for the actual collapse, including its integral index.

For a fully marked certificate we represent a base cycle $b_J$ by
Poincar\'e duality from the Alexander--Whitney cup cocycle on the
complementary coordinates.  On an ordered $p$-simplex its coefficient
is
\[
\operatorname{sgn}(J,J^c)
 \prod_{r=1}^p(v_r-v_{r-1})_{(J^c)_r}.
\]
Tensoring with an angular exterior vector and applying (1.11.1)
gives the matrices exported by \texttt{cusp\_marked\_collapse.py}.
They annihilate $\bigwedge^kM_0-I$ exactly.  The actual collapse also
annihilates that operator, since turning the smoothing parameter
does not change its image in the central fibre.  The computed ranks
equal the ranks of the monodromy coinvariants; primitivity therefore
identifies both kernels with the saturation of its image.
Consequently the exported matrices are matrices of the actual
collapse in suitable integral target bases.  This last argument also
eliminates any ambiguity from choosing splittings of the filtration.
\end{proof}

\begin{statement}{P1.12}{Proved}
The fourfold $M_c$ has torsion-free integral homology
\[
H_k(M_c;\mathbb Z)=
\begin{cases}
\mathbb Z,&k=0,8,\\
\mathbb Z^4,&k=4,\\
0,&k=1,2,3,5,6,7.
\end{cases}
\]
In particular it is three-connected.
\end{statement}
\begin{proof}
We compute the complement $Y$ of the open cusp cap, then attach the
cusp using only the degrees needed to prove three-connectivity.
This avoids assigning a value to a boundary suspension that is not
determined by the fibre specialization alone.

Write $L_k=\bigwedge^k\Lambda$ and $W_{i,k}=\bigwedge^kM_i$.
The cellular model of the main torus bundle over a graph with two
oriented edges has
\[
G_k=L_k\oplus L_{k-1}\oplus L_{k-1},\qquad
d_G(f,x_1,x_2)=((W_{1,k-1}-I)x_1+(W_{2,k-1}-I)x_2,0,0).
\tag{1.12.1}
\]
The boundary mapping torus at $i$ has
$B_{i,k}=L_k\oplus L_{k-1}$ and
\[
d_{B_i}(f,x)=((W_{i,k-1}-I)x,0).
\tag{1.12.2}
\]
The finite boundary maps to $G$ by identity on $L_k$ and inclusion
of the suspension into its respective edge.  The cusp boundary,
corresponding to $(g_1g_2)^{-1}$, maps by
\[
j_{0,G}(f,x)=(f,-W_{2,k-1}W_{0,k-1}x,-W_{0,k-1}x).
\tag{1.12.3}
\]
This follows by traversing the second edge backwards and then the
first backwards, transporting the fibre chain to each initial end.
Indeed its boundary is $W_{0,k-1}-I$ because
\[
-(W_1-I)W_2W_0-(W_2-I)W_0=W_0-I.
\]

Here are full models for the finite pieces and their attaching maps.
Put $\varepsilon_i=\ell(v_i)$ and
$P_i=(\varepsilon_iv_i,e_1,\ldots,e_5)$.
This is unimodular.  It splits
$\Lambda=\mathbb Z(\varepsilon_iv_i)\oplus K$, $K=\ker\ell$,
and $P_i^{-1}M_iP_i=1\oplus R_i$.
The cap retracts onto the mapping torus of $R_i$ on $T_K$.
Its chains are
\[
E_{i,k}=\bigwedge^kK\oplus\bigwedge^{k-1}K,
\qquad d_{E_i}(f,x)=((\bigwedge^{k-1}R_i-I)x,0).
\tag{1.12.4}
\]
In these split coordinates, with $v=\varepsilon_iv_i$, the finite
boundary-to-cap chain map is specified on all summands by
\[
\begin{array}{c|c}
\text{boundary chain}&\text{cap chain}\\\hline
x\in\bigwedge^kK&(x,0)\\
v\wedge x&(0,\varepsilon_i\sum_{j=0}^{m_i-1}
                  (\bigwedge^{k-1}R_i)^jx)\\
g_i\wedge x&(0,x)\\
g_i\wedge v\wedge x&0.
\end{array}
\tag{1.12.5}
\]
Changing back uses
$\operatorname{diag}(\bigwedge^kP_i^{-1},
\bigwedge^{k-1}P_i^{-1})$.
Geometrically this map is
$(t,y,\theta)\mapsto[y,\theta+\varepsilon_im_it]$ into the mapping
torus.  The $t$-circle sweeps $m_i$ strips, giving the norm sum; the
$g_i$-circle sweeps one strip; their product maps to a single
base direction, giving zero on the last row.  This proves the map
including its sign and the possible finite-cap torsion.

We justify that these are integral cellular models, rather than a
formal use of invariant rational forms.  In the bases with columns
\[
K_1=\begin{pmatrix}
-1&-1&0&-1&0\\0&1&0&1&1\\0&1&0&0&0\\
0&-1&-1&0&-1\\0&0&1&0&0
\end{pmatrix},\qquad
K_2=\begin{pmatrix}
-1&0&-1&0&-1\\-1&1&0&1&-1\\0&1&-1&1&-1\\
0&-1&0&-1&1\\1&-1&1&0&1
\end{pmatrix},
\]
both unimodular, $R_1$ is the permutation of five coordinates in a
cycle, and $R_2$ is a four-cycle together with multiplication by $-1$
on the fifth coordinate.  The minimal product cell structure of
circles is preserved by these signed permutations; its cellular
boundary is zero and its action on cells is exactly the exterior
power.  Formula (1.12.5) is obtained by subdividing its swept
base interval into the indicated strips.  The common torus can be
marked into these two structures by based cellular approximations.
Their cellular maps are the exterior-power matrices because the
cellular differentials vanish.  Based maps between these aspherical
tori with a fixed map on $\pi_1$ have a contractible space of
homotopies: lifting to the universal covers gives equivariant maps
to a contractible vector space and linear interpolation relative to
the base point.  Thus the choices introduce no extra attaching class
on the graph.  Composing the two edge cylinders gives exactly
(1.12.3).  The explicit conjugators were found in the bounded search
$[-1,1]^5$ and then verified by exact multiplication in
\texttt{finite\_permutation\_models.py}.

Let $F_k:B_{1,k}\oplus B_{2,k}\to G_k\oplus E_{1,k}\oplus E_{2,k}$
have the main boundary maps as its first row, and the negatives of
(1.12.5) in the corresponding two cap rows.  The chain complex of
$Y$ is the mapping cone
\[
C_k(Y)=G_k\oplus E_{1,k}\oplus E_{2,k}
       \oplus B_{1,k-1}\oplus B_{2,k-1},\qquad
d_Y=\begin{pmatrix}d_G\oplus d_{E_1}\oplus d_{E_2}&F_{k-1}\\
0&-d_{B_1}\oplus-d_{B_2}\end{pmatrix}.
\tag{1.12.6}
\]
Equations (1.12.1)--(1.12.6), minors for exterior powers, and the input
matrices specify every entry of every differential.
The full arrays are exported to \texttt{complement\_export.json}; its
program verifies $d^2=0$ in all degrees and computes Smith
decompositions.  Chain ranks in degrees zero through eight are
$(3,22,71,132,155,118,57,16,2)$.
The resulting homology of $Y$ is free, with ranks
$(1,1,3,9,12,9,3,0,0)$.

The boundary $B_0$ has
\[
H_1(B_0)=\mathbb Z^4,\qquad H_2(B_0)=\mathbb Z^{12},\qquad
H_3(B_0)=\mathbb Z^{18}\oplus\mathbb Z/2.
\]
The homomorphism $j_Y:H_k(B_0)\to H_k(Y)$ from (1.12.3) is
surjective for $k=1,2,3$.  In degrees two and three, the program
\texttt{complement\_kernel\_check.py} supplies integral cycle bases
of its kernel whose entire suspension block in (1.12.2) is zero.
This is an exact zero matrix, not a vanishing after tensoring with
$\mathbb Q$.  Therefore these kernels are represented by cycles of
the nearby torus itself.  The restriction of the actual map to the
cusp cap on these kernels is consequently the already proved
collapse of P1.11, regardless of its value on other boundary
suspension cycles.

For transparency, in the Smith bases exported with these cycles the
combined map in degree one is
\[
\begin{pmatrix}20&1&0&0\\1&0&-1&-1\\1&0&0&0\\0&0&2&1\end{pmatrix},
\qquad\det=1.
\tag{1.12.7}
\]
The cusp meridian maps to zero here because the zero section extends.
In degree two $\ker j_Y=\mathbb Z^9$ and its map to $H_2(C)$ is
\[
\begin{pmatrix}
0&0&1&-1&0&-1&0&0&0\\
0&0&1&0&0&-1&0&1&0\\
0&0&1&-1&0&-2&0&1&0\\
0&0&0&0&-1&-1&0&0&-1\\
1&1&-1&0&0&1&1&0&0\\
1&2&0&0&1&0&0&0&0\\
-1&-1&0&0&0&0&0&0&0\\
0&0&0&0&1&0&0&0&0\\
0&0&0&0&0&-1&0&0&0
\end{pmatrix},\qquad\det=1.
\tag{1.12.8}
\]
In degree three $\ker j_Y=\mathbb Z/2\oplus\mathbb Z^9$.
Its collapse matrix, with the torsion generator first, is
\[
\begin{pmatrix}
0&1&0&1&0&0&0&0&0&1\\
0&1&-1&0&0&0&-1&0&-1&0\\
0&0&0&0&0&0&1&0&1&0\\
0&0&0&-1&1&-1&0&-1&0&0\\
0&0&0&0&0&0&0&1&0&0\\
0&0&0&0&0&0&1&0&0&0\\
0&0&0&0&0&1&0&1&0&0\\
0&0&1&0&1&-2&1&-1&2&0\\
0&1&0&-1&1&-2&0&-1&0&0
\end{pmatrix}.
\tag{1.12.9}
\]
The determinant of its last nine columns is one.  In particular the
torsion of the boundary and the second finite cap has been included
throughout the calculation.

We now apply Mayer--Vietoris to $M_c=Y\cup C_{\rm cap}$.
Because $j_Y$ is onto, subtracting a preimage of the $Y$-coordinate
identifies the cokernel of $(j_Y,j_C)$ with
$H_k(C)/j_C(\ker j_Y)$; its kernel is
$\ker j_Y\cap\ker j_C$.
Equations (1.12.7)--(1.12.9) therefore say that the combined map is an
isomorphism in degrees one and two and is onto in degree three.
The exact sequence yields $H_2(M_c)=H_3(M_c)=0$.
We already proved $\pi_1(M_c)=0$, and successive applications of
Hurewicz now give $\pi_2(M_c)=\pi_3(M_c)=0$.

Integral Poincar\'e duality and the universal coefficient theorem
give $H_5=H_6=H_7=0$.  They also show that $H_4$ is free: its possible
torsion would be $\operatorname{Ext}(H_4,\mathbb Z)$ in $H^5$, which
vanishes by duality with $H_3=0$.
Finally $\chi=6$ gives $\operatorname{rank}H_4=4$.
This proves all claimed integral homology groups without making an
assumption about the higher-degree cusp suspension chain term.
\end{proof}

\paragraph{The retained diagnostic complex.}
The earlier file \texttt{global\_candidate\_complex.py} sets the entire
cusp suspension map to zero and has the same free Betti numbers.
That file and its output are retained as a diagnostic record.  They
are not used in P1.12 to identify the actual higher-degree gluing.
P1.19 supplies the requested independent full global cell computation.
It proves that the local secondary cusp term is nonzero, gives the
explicit integral reduction to the corrected global matrix, and
recovers every homology group without global Poincar\'e duality.

\begin{statement}{P1.13}{Proved}
With the complex orientation the fourfold has
\[
\sigma(M_c)=2,\qquad H^4(M_c;\mathbb Z)\cong I_{3,1},\qquad
w_2(M_c)=0,
\]
and there is a diagonal integral basis in which
$p_1/2=(1,1,1,1)$.
Its Chern numbers are
$c_4=6$, $c_2^2=2$, and
$c_1^4=c_1^2c_2=c_1c_3=0$; its Pontryagin numbers are
$p_1^2=8$, $p_2=14$.
\end{statement}
\begin{proof}
P1.12 implies $c_1=c_3=0$ as integral cohomology classes;
$c_4[M_c]=\chi(M_c)=6$.  The Chern-root identities give
$p_1=-2c_2$, $p_2=c_2^2+2c_4$.
Hirzebruch, \emph{Topological Methods in Algebraic Geometry},
Theorem 24.5.4 and its corollary, give integrality of the Todd genus
for this compact almost complex manifold.  Expanding
$\prod x_j/(1-e^{-x_j})$ to degree four gives
\[
\frac{3c_2^2-6}{720}\in\mathbb Z,
\quad\text{hence}\quad c_2^2=2+240k\quad(k\in\mathbb Z).
\]
The signature theorem in the same source, Theorem 8.2.2, applied to
this closed oriented smooth eight-manifold gives
\[
\sigma=\frac{7p_2-p_1^2}{45}
      =\frac{3c_2^2+84}{45}=2+16k.
\]
The rank-four middle form has $|\sigma|\leq4$, so $k=0$.
This proves the displayed characteristic numbers.

Duality makes the middle form unimodular.  Its signature is two and
its positive and negative indices are three and one.  An even
unimodular form has signature divisible by eight, by Wall,
\emph{Classification of $(n-1)$-connected $2n$-manifolds},
Proposition 3, p.~173.  Hence it is odd.  Proposition 4 on that page
classifies indefinite unimodular forms by rank, signature, and
parity, and gives $I_{3,1}$.
Since $H^2(M_c;\mathbb Z/2)=0$, its $w_2$ is zero.
For the complex tangent bundle,
$w_4=c_2\bmod2=(p_1/2)\bmod2$.
The Wu identity in degree four, with $w_1=w_2=0$, says that $w_4$ is
characteristic for the middle pairing.  Thus the vector $p_1/2$ is
characteristic and has square $2$.

We give the required integral isometry of this pair explicitly as
a terminating procedure.  In a diagonal basis a characteristic vector
has four odd coordinates.  Change their signs to write it as
$(a,b,c,d)$ with all positive and
$a^2+b^2+c^2=d^2+2$.
For $d>1$,
\[
d<a+b+c<3d.
\]
The first inequality follows by squaring and using the positive cross
terms; the second follows from Cauchy--Schwarz and $d>1$.
Reflection in the integral norm-two vector $(1,1,1,1)$ sends this
vector to itself minus $(a+b+c-d)(1,1,1,1)$.
Its fourth coordinate is $2d-a-b-c$, whose absolute value is strictly
smaller than $d$.  The reflection preserves the integral lattice and
characteristic parity.  Repeating, with sign changes at each step,
terminates at $d=1$, when the norm equation forces $a=b=c=1$.
This supplies an integral isometry taking $p_1/2$ to $(1,1,1,1)$,
and proves the claimed pair, not just its square.
\end{proof}

\begin{statement}{P1.14}{Known}
The constructed fourfold is orientation-preservingly diffeomorphic to
$3\mathbb HP^2\#(-\mathbb HP^2)$.
\end{statement}
\begin{proof}
We use Wall's almost-diffeomorphism classification in the numbered
form of Conway--Crowley--Powell--Sixt,
\emph{Simply connected manifolds with large homotopy stable classes},
Theorem 3.7(1), with its Definition 3.6.  It identifies closed oriented
$(4m-1)$-connected $8m$-manifolds up to an orientation-preserving
diffeomorphism after connected sum with a homotopy sphere by their
intersection form together with the stable normal homomorphism.
We read the statement and its reduction to Wall's theorem at
\url{https://www.maths.gla.ac.uk/~mpowell/large-homotopy-stable-class-v3.pdf},
pp.~12--14, and read the original almost-closed classification on
Wall's p.~170.

Here $m=1$.  Both manifolds are closed, oriented, smooth and
three-connected, the first by P1.12 and the second by its standard
connected-sum cell structure.  A class in $H_4$ can be represented by
an embedded four-sphere as in the theorem.  Its stable normal bundle
is the restriction of the ambient tangent bundle, since the sphere
is stably parallelizable.  The group of stable oriented bundles on
$S^4$ is $\pi_3(SO)=\mathbb Z$, detected by $p_1/2$ after the same
generator convention is chosen on both manifolds.  Thus the extended
symmetric form in this dimension is exactly the pair computed in
P1.13.

For $\mathbb HP^2$, take the generator $\alpha$ with
$p_1=2\alpha$.  Its square evaluates to one and $p_2=7\alpha^2$;
these formulas follow by expanding
$(1+\alpha)^6/(1+4\alpha)$ through degree eight, and are also verified
in Mandal, arXiv:2609.27929, Lemma 2.1.
The connected sum has form $I_{3,1}$ and
$p_1/2=(1,1,1,1)$, since the Pontryagin class restricts to the given
class on each punctured summand and the last summand carries the
reversed fundamental orientation.  The theorem's two extended forms
are therefore isomorphic.  We obtain an orientation-preserving
diffeomorphism to the target connected-summed with a homotopy
eight-sphere.  P1.6 absorbs that sphere by an actual oriented
diffeomorphism.  This proves the asserted diffeomorphism type with
no remaining ambiguity.
\end{proof}

\begin{statement}{P1.15}{Proved}
The rank-six flat extension, as a complex flat bundle with its invariant
line, has the following explicit algebraic connection model.  Set
$\omega_k=x^kdx/y$, $0\leq k\leq3$, and use the quotient frame
\[
(w_2,w_1,w_0,w_{-1},w_{-2})
 =(\omega_{01},\omega_{02},\omega_{03},\omega_{13},\omega_{23}),
\qquad \omega_{12}=-3\omega_{03}.
\]
These are meromorphic frames of
$(L^2,L,\mathcal O,L^{-1},L^{-2})$, where $L=T^2$ on the root
orbifold of orders $5,4$.  If a frame is written as a row of vectors,
its derivative is right multiplication by
\[
\Gamma_6=\begin{pmatrix}0&b\\0&\Gamma_W\end{pmatrix},\qquad
b=\frac1{\lambda-\lambda_*}
 \left(-\frac{25}{6\lambda^2},-\frac{125}{48\lambda},0,1,0\right),
\tag{1.15.1}
\]
where
\[
\resizebox{\textwidth}{!}{$\displaystyle
\Gamma_W=\frac1{3125\lambda+256}
\begin{pmatrix}
-8(625\lambda+48)/\lambda&-8&10\lambda&75\lambda^2/2&0\\
-125/(2\lambda)&-4(625\lambda+48)/\lambda&8&0&75\lambda^2/2\\
300/\lambda^2&-375/\lambda&0&-48&60\lambda\\
120/\lambda^3&0&125/(2\lambda)&4(625\lambda+48)/\lambda&-8\\
0&120/\lambda^3&50/\lambda^2&-125/(2\lambda)&8(625\lambda+48)/\lambda
\end{pmatrix}.$}
\tag{1.15.2}
\]
This model specifies a canonical algebraic splitting, denoted
$\sigma_b$.  The compatible contraction used for the real parameter is
$\alpha=(625\pi i/128)\alpha_{\rm alg}$, where
$\alpha_{\rm alg}(w_0)=4/3$ and it vanishes on $w_1,w_2$.
\end{statement}
\begin{proof}
Reduction modulo exact differentials is completely explicit.  To
differentiate $\omega_k$, solve the polynomial identity
\[
-\tfrac12x^k\partial_\lambda f
 =f\sum_{j=0}^3G_{kj}x^j+R_k'f-\tfrac12R_kf',
 \qquad\deg R_k\leq4,
\]
where $f=x^5+\lambda^3x+\lambda^4$.  The last two terms are the
numerator of $d(R_k/y)$.  Matching the nine coefficients gives the
four-by-four Gauss--Manin matrix exactly.
At the branch point at infinity, a local coordinate $z=x^{-1/2}$
gives $\omega_k=-2z^{2-2k}(1+O(z^8))\,dz$.
Residues of a primitive times a differential give the alternating
pairings $J_{03}=4/3$, $J_{12}=4$, with their negatives in reverse
order and all other entries zero.  This form is nondegenerate.
The invariant bivector is therefore $3\omega_{03}+\omega_{12}$.
Taking the exterior-square connection and this quotient gives
(1.15.2), as checked symbolically in \texttt{de\_rham\_connection.py}.

The root-chart powers computed from $x^kdx/y$ are
$u^{6k-9}$ at zero and $u^{6-4k}$ at infinity.
Their nondegenerate pairing shows that the four resulting regular
frames form the full de Rham bundle on those charts.  At the node
the computed connection has a simple pole and square-zero residue;
hence these are the nilpotent logarithmic extension frames there.
Thus the four line bundles are
$T^3,T,T^{-1},T^{-3}$, and their quotient exterior-square lines are
the stated powers of $L$.

All $H^1(L^j)$ for $j=0,\pm1,\pm2$ vanish by the coarse-degree
calculation already used in P1.8.  Consequently every holomorphic
extension of this $W$ by $\mathcal O$ splits as a bundle.
In such a splitting its flat connection is an upper block extension
by a row $b$.  Regularity at the finite root charts and a logarithmic
pole at the cusp require, entry by entry,
\[
b=\frac1{\lambda-\lambda_*}
 (a\lambda^{-2},b'\lambda^{-1},0,c,d\lambda).
\]
For example the $L^j$ entry, after $\lambda=u^4$, has exponent
$3j+3+4p$ for a term $\lambda^p/ (\lambda-\lambda_*)$;
the infinity chart requires $-4j-1-5p\geq0$.
These two inequalities leave exactly the four displayed terms.
The equation that the full cusp residue have square zero is
$b_{\rm res}(\Gamma_W)_{\rm res}=0$.
Its solution space has basis
\[
(-25/6,-125/48,1,0),\qquad(5/3,-25/3,0,1).
\]
The sole nonconstant holomorphic splitting change is the row selecting
$w_0$, and its connection coboundary is the third row of $\Gamma_W$.
In these two basis vectors that row is $(-48b_1+60b_2)/3125$.
Thus there is one nontrivial flat extension direction modulo splitting
changes and rescaling of the invariant line.  The first vector has
rank-three square-zero cusp residue, so it represents that direction.
The supplied extension is nontrivial, since its dual has no common
invariant vector by P1.1.  It is therefore isomorphic to (1.15.1).
This proves the connection identification without inferring it from
the subsequent numerical matching.

For the last normalization, the kernel of the quotient cusp residue
on $(w_2,w_1,w_0)$ is generated by
$v=(-4096/1171875,64/1875,1)^t$.
It satisfies $b_{\rm res}(v)=625/192$ and
$\alpha_{\rm alg}(v)=4/3$.
The logarithmic monodromy is $2\pi i$ times this residue.
The exact cusp identity of P1.8 characterizes the contraction on this
line by $N_V(\widetilde v)=\alpha(v)e$.
Since the space of global contraction forms is one-dimensional, the
ratio $(625/192)/(4/3)=625/256$ fixes its scalar everywhere and gives
the displayed $625\pi i/128$.
\end{proof}

\paragraph{An explicit formula for $h$ and the normalization comparison.}
Let $Y'=Y\Gamma_6$.  Normalize $Y(1)=X$ by requiring that continuation
around the specified meridians intertwine exactly with $A,B$, and that
$Xe_0=e_0$.  Existence follows from P1.15; uniqueness follows from the
scalar centralizer, whose exact linear equations are in
\texttt{centralizer\_torsion.py}.
At each nonspecial point there is a unique pair $v_2,v_1\in\mathbb C$
such that the lower five coordinates of
$Y(0,v_2,v_1,-i,0,0)^t$ are real.
The intersection-line proof in P1.8 and the normalization of $\alpha$
prove existence and uniqueness.  Therefore the following is a fully
specified period formula for the real function associated with
$\sigma_b$:
\[
h_b(\lambda)=\frac{96}{625\pi}
 \operatorname{Im}\bigl(Y_{01}v_2+Y_{02}v_1-iY_{03}\bigr).
\tag{1.15.3}
\]
Its periods for the real parameter $c$ are the three columns
$Y_{*,1},Y_{*,2},Y_{*,3}-(625\pi/96)c\,e_0$.

This is a different splitting normalization from the one chosen in
P1.8.  To compare them exactly, parallel-transport the fixed marking
from $\lambda=1$ to $\lambda=\lambda_*/2$ along a specified path in the
upper half-plane, and put
$a=(\sigma_b(g))_e$ at that point, where
$g=(96/(625\pi i))w_0$ and $\alpha(g)=1$.
Then the splitting of P1.8 is
$\sigma=\sigma_b-a\alpha e$, so $h=h_b-\operatorname{Im}a$.
For the full complex splitting parameter $\xi$, the comparison is
$\xi_b=\xi-a$.  This gives the actual translation between the two
families, including its possible real part; it does not silently
identify their real parameter slices.

\paragraph{Numerical record at the first checkpoint.}
The code \texttt{pencil\_period\_numerics.py} independently evaluates
the four branch-tree periods of $dx/y$ and $x\,dx/y$ at $\lambda=1,2$
by endpoint-resolved tanh--sinh quadrature and by convergent binomial
series with exact cosine moments, at ninety decimal digits.
The lifted connection computation in
\texttt{connection\_monodromy.py} gives finite-order residuals below
$3\cdot10^{-96}$ and intertwiner residual below $8\cdot10^{-97}$.
Its resulting lifted periods and determinants are exported in
\texttt{lift\_period\_values.json}.  In the $\sigma_b$ normalization,
\[
\begin{split}
h_b(1)&=-1.03427662241632179949424884537834605259\ldots,\\
h_b(2)&=-1.04792484985922734646718835593486351800\ldots.
\end{split}
\]
The preliminary loop calculation uses midpoint recentering and residual
checks.  P1.16 supplies the rigorous interval continuation and the
independent Frobenius computation of the lifted rank-six periods;
the curve periods have the additional independent quadrature and
cosine-moment evaluations just described.

\begin{statement}{P1.16}{Proved}
In the canonical splitting $\sigma_b$ of P1.15,
\[
\begin{aligned}
-1.0342766224163217994942488453783461&<h_b(1),\\
h_b(1)&<-1.0342766224163217994942488453783460.
\end{aligned}
\]
In particular $h_b(1)<-1$.  The parameter $c=5/(3\pi)$ is outside the
global transversality domain.  The lifted period matrices at
$\lambda=1,2$ have been evaluated by two different continuation methods
agreeing to more than thirty decimal digits.
\end{statement}
\begin{proof}
The file \texttt{certified\_lift\_bound.py} propagates complex balls
without dropping their radii.  We describe the explicit truncation
bound so the numerical inequality has a mathematical certificate.
On a Taylor step from $z$ to $z+h$, choose a disk of radius $R$ avoiding
the two finite poles.  Use a fixed positive diagonal matrix $S$ on
this disk.  Each entry of the connection is a polynomial divided by
$\lambda^p(\lambda-\lambda_*)$; the triangle inequality gives a
rigorous bound $M$ for the maximum row-sum norm of
$S^{-1}\Gamma_6S$ on the disk.  Along a radius, the integral equation
and Gronwall give a bound $\exp(MR)$ for the fundamental solution
normalized to the identity at its center.  Cauchy's coefficient
estimate then bounds the omitted Taylor tail after degree $N$ by
\[
\exp(MR)\frac{(|h|/R)^{N+1}}{1-|h|/R}.
\tag{1.16.1}
\]
The program adds this error to each entry with factor $S_i/S_j$.
All coefficient evaluations, matrix products, and linear solves retain
their outward-rounded complex-ball enclosures.

For both finite meridians we use rational polygonal paths, and choose
the order adaptively until (1.16.1) is below $10^{-65}$.
There are $152$ steps in total and the largest Taylor order is $239$.
The thirty independent equations determining the normalized
intertwiner are solved by interval matrix inversion; exact existence
and uniqueness are P1.15 and the scalar-centralizer calculation.
The real intersection vector in (1.15.3) is likewise obtained by an
invertible four-equation subsystem, whose exact consistency follows
from P1.8.  Substitution into (1.15.3) gives the certified enclosure
\[
h_b(1)\in
[-1.03427662241632179949424884537834605259496477790218821132
 \ \mathbin{\pm}\ 8.11\cdot10^{-57}],
\]
which implies the stated slightly wider decimal bounds.
Every step radius, norm bound, order and tail bound is exported in
\texttt{certified\_lift\_bound.json}.  Since $5/(3\pi)<1$, the value
$h_b(1)+5/(3\pi)$ is negative, whereas $h_b+c$ tends to positive
infinity at the cusp.  A path joining these points must contain a
zero.  This proves inadmissibility, not merely failure of a numerical
test at one point.

For the independent period evaluation,
\texttt{frobenius\_period\_check.py} uses regular finite-root frames
and their deck eigencharacters, without integrating a closed loop.
At zero the frame powers are $(0,6,3,0,-3,-6)$ in $u$ with
$\lambda=u^4$, and at infinity they are
$(0,-8,-4,0,4,8)$ with $\lambda=u^{-5}$.
The transformed connections are regular at $u=0$ and have denominators
$3125u^4+256$ and $3125+256u^5$, respectively.
Their Taylor recursions therefore give the connection matrices from
the center to $u=1/4$ and $u=1/2$.  Conjugating the exact diagonal
deck characters by these matrices gives the finite monodromies.
Their normalized intertwiner agrees with the closed-loop calculation
to $4.5\cdot10^{-94}$ in maximum entry norm.
The resulting period planes at $\lambda=1,2$ and their real
transversality determinants are all recorded in the accompanying
JSON files; for example at $c=10$ the determinants are purely imaginary
with positive imaginary parts
\[
\begin{gathered}
88279007523.2167451846066607088299709479845\ldots,\\
3178499608.3637608517583236010752861132441\ldots.
\end{gathered}
\]
These last two local checks do not by themselves prove that $c=10$
is globally admissible.  Global admissibility is exactly P1.8's
inequality involving the minimum.  The independent root-chart
calculation also gives
\[
\begin{split}
h_b(\infty)&=-1.094540250311402780116268263054945703725533\ldots,\\
h_b(0)&=-0.931793376356523505430795105796011114990828\ldots.
\end{split}
\]
P1.18 proves, by a complete interval covering argument, that the
first value is the unique global minimum.
\end{proof}

\begin{statement}{P1.17}{Proved}
For the cusp marking $S$ fixed in P1.1 and the corresponding Kuhn fan,
none of the completed fourfolds in this family admits an
anti-holomorphic automorphism.  In particular none admits a real
structure.  More strongly, two members with this same cusp marking
are never anti-biholomorphic.
\end{statement}
\begin{proof}
We first recover the base intrinsically.  An effective divisor on
$M_c$ cannot be horizontal: its first Chern class vanishes by
$H^2(M_c;\mathbb Z)=0$, whereas its nonzero effective restriction to a
general torus fibre would have positive pairing with the square of a
flat K\"ahler class.  The pole divisor of any meromorphic function is
therefore vertical.  On every regular fibre outside that divisor the
function is holomorphic and hence constant.  It factors through a
meromorphic function on the base, which is rational on
$\mathbb P^1$.  Thus the function field is exactly $\mathbb C(\lambda)$.
An anti-biholomorphism consequently induces an anti-M\"obius map of
the base.  The two distinct multiplicities $5,4$ and the singular cusp
support distinguish the three special values, so this base map fixes
$\infty,0,\lambda_*$.  As those points are real, the map is precisely
$\lambda\mapsto\bar\lambda$.

A holomorphic map between compact complex tori is affine: after
lifting to their vector spaces, its derivative is periodic and
therefore constant, since it descends to holomorphic functions on a
compact torus.  Applying this to an anti-holomorphic map after
conjugating its source gives the same assertion for its anti-linear
lift.  Its induced integral map on $H_1$ is constant along the
regular base.  Complex conjugation reverses the two finite meridians,
so this matrix must simultaneously conjugate $M_i$ to $M_i^{-1}$.
Solving these exact integral linear equations gives only $R$ and
$-R$ among unimodular matrices, where
\[
R=\begin{pmatrix}
1&0&0&0&0&0\\
2&0&0&0&-1&0\\
-1&0&-1&0&0&0\\
-1&0&0&-1&0&0\\
2&-1&0&0&0&0\\
-1&1&1&1&1&1
\end{pmatrix},\qquad R^2=I.
\]
Indeed the simultaneous intertwiner space over $\mathbb Q$ has
dimension one and this primitive matrix is unimodular; determinant
one in absolute value permits only the scalar factors $\pm1$.
The equations and this certificate are in
\texttt{centralizer\_torsion.py} and
\texttt{real\_structure\_check.py}.

The finite affine cap excludes $G=R$.  Normalize the pullback by
$t=u^{m_i}$; any base-preserving isomorphism lifts to this canonical
root cover and conjugates its deck generator to the inverse.
At its central torus the translation part $t_i$ would satisfy
\[
(I-M_i^{-1})t_i=-(v_i+Gv_i)/m_i\pmod\Lambda.
\]
Applying $\ell$, and using $\ell R=\ell$, gives
$0=-2\ell(v_i)/m_i$ in $\mathbb R/\mathbb Z$, impossible already
for $m_i=5$.  Hence $G=-R$ is necessary.

We now use the cusp marking, retaining the conjugation of based
meridians.  Conjugation sends
$g_0=(g_1g_2)^{-1}$ to $g_2g_1=g_2g_0^{-1}g_2^{-1}$.
Thus the matrix in a fixed local cusp frame is
$H=M_2^{-1}G$, and it anti-conjugates $M_0$.
Exact multiplication in the marking $S$ gives
\[
S^{-1}HS=
\begin{pmatrix}
1&0&0&-1&-1&0\\
0&1&0&-2&0&0\\
1&-1&-1&1&-1&-1\\
0&0&0&-1&0&0\\
0&0&0&0&-1&0\\
0&0&0&-1&1&1
\end{pmatrix}
=\begin{pmatrix}P&Q\\0&-P\end{pmatrix},\quad
P=\begin{pmatrix}1&0&0\\0&1&0\\1&-1&-1\end{pmatrix}.
\tag{1.17.1}
\]
Changing the based cusp path by a cusp loop multiplies by a power of
$M_0$ and does not change $P$.

We explain why an arbitrary translation cannot remove the obstruction
in (1.17.1).  In the additive torus coordinates whose first periods
are the $a_j$, the anti-linear map must have the form
$\zeta' =P\bar\zeta+t(q)$, because its values on these three real
periods already determine its linear part.  In multiplicative
coordinates it is
$z'=u(q)\bar z^{-P}$.
The source zero section limits to a smooth point of the central fibre.
An isomorphism takes it to another smooth point.  The central fibre
is irreducible, and its smooth locus is exactly the dense torus of
its toric normalization: every toric boundary divisor is a
normal-crossings self-identification.  In a chart about this image,
the translation therefore has the form $q^v$ times nonvanishing
holomorphic units, with $v\in\mathbb Z^3$.  Translating by the period
$v$ changes the selected component chart; the units act by torus
translation.  Neither changes the linear action $-P$ on the fan.

The primitive rays of the normal fan of that component are
\[
\{\pm v:0\ne v\in\{0,1\}^3\}.
\]
An extending torus monomial isomorphism must permute these rays.
This can be seen directly on a one-parameter arc $z=t^v$ meeting the
generic point of its boundary divisor: if the image vector lies in
the interior of a cone of dimension at least two, that divisor is
contracted, which is impossible for an isomorphism.  Complex
conjugation itself preserves the toric charts, so the same criterion
applies to the anti-monomial map above.
But $-Pe_2=(0,-1,1)^t$ is not one of the displayed rays.
This is a contradiction.  We have exhausted both possible integral
linear parts, so no anti-biholomorphism exists.
The assertion is explicitly about the chosen $S$ and its fan; it does
not identify different toroidal compactifications without a map
between them.
\end{proof}

\begin{statement}{P1.18}{Proved}
In the algebraic splitting $\sigma_b$, the function $h_b$ has its
unique global minimum at the finite-root point $\lambda=\infty$.
The minimum is the exact Frobenius period constant
\[
 h_{b,\min}=-\frac{96}{625\pi}\operatorname{Re}(XH_\infty)_{03},
 \qquad
 h_{b,\min}=-1.0945402503114027801162682630549457037\ldots,
\tag{1.18.1}
\]
where $X$ is the normalized integral intertwiner of P1.15 and
$H_\infty$ is the regular infinity-root limit of its fundamental
solution, normalized to be the identity at $\lambda=1$ before
multiplication by $X$.  Formula (1.18.1), the rational connection,
and those two normalizations specify the constant exactly.
A certified enclosure has radius $4.04\cdot10^{-56}$ about the
56-digit midpoint retained in
\path{code/part1/minimum_certificate_summary.json}.
Consequently the exact domain of completed transverse periods is
\[
 c>-h_{b,\min},\qquad
 \operatorname{Im}\xi>-h_{b,\min}
\]
for the real and complex parameters, respectively.
The member $c=10$ used in the numerical period checks is therefore
admissible on the whole completed base.

For the point-normalized splitting of P1.8, continued from
$1$ to $\lambda_*/2$ by the polygon
$1,1+i,\lambda_*/2+i,\lambda_*/2$, the exact comparison constant
$a_{\rm split}$ of (1.15.4) has certified value
\[
\begin{split}
 \operatorname{Re}a_{\rm split}
   &=-0.7099885679126427771723473455012308149\ldots,\\
 \operatorname{Im}a_{\rm split}
   &=-0.9134017059364153362143229223828373269\ldots.
\end{split}
\]
Thus the minimum in that splitting is
$-0.1811385443749874439019453406721083768\ldots$,
attained at the same base point.  Its real parameter interval is
\[
 c>0.1811385443749874439019453406721083768\ldots.
\]
The full parameter comparison remains
$\xi_b=\xi_{\rm original}-a_{\rm split}$; both its real and imaginary
parts are necessary to identify the two normalized families.
\end{statement}
\begin{proof}
We give the error bounds and the complete covering argument behind
the interval computation.  This is a finite proof of the minimum,
not an inference from the preliminary grid search.  All coefficients,
accepted rectangles and local continuation bounds are retained in the
four \texttt{certified\_minimum\_*} programs and their JSON outputs.
The file \texttt{minimum\_certificate\_summary.json} records their
SHA-256 digests as well as sharp interval endpoints.

Put $a=|\lambda_*|$.  The normalized root coordinates are
\[
 z_\infty=(a/\lambda)^{1/5},\qquad
 z_0=(\lambda/a)^{1/4}.
\]
In the regular root frames of P1.16 their connections have the form
\[
 G(z)=\frac{\sum_{j=0}^{m-1}B_jz^j}{1+z^m},\qquad m=5,4.
\tag{1.18.2}
\]
The entries of $B_j$ are the exact algebraic numbers obtained by
substituting the positive real number $a^{\pm1/m}$ in the rational
connection.  The generator of the finite deck group acts on these
frames by their stated eigencharacters.  Since its flat matrix is
real and fixes $e$, it leaves $h_b$ invariant.  It therefore suffices
to cover the sector $0\leq\arg z\leq2\pi/m$ in each root disc.

Use the positive diagonal weights
\[
 (1,1/128,1/128,1/32,1/64,1/32)
\]
and the corresponding weighted maximum row norm.  Exact ball
arithmetic gives
\[
 \left\|\sum_j |B_j|\right\|<1.447\quad(m=5),\qquad
 \left\|\sum_j |B_j|\right\|<1.135\quad(m=4).
\]
Here the absolute values are entrywise, before taking the row norm.
Let $B$ denote the actual enclosed bound.  The scalar inequality
\[
 \frac1{1-t^m}\leq\frac1{m(1-t)}+\frac{m-1}{m},\qquad0\leq t<1,
\]
follows on multiplying by $m(1-t)(1+\cdots+t^{m-1})$;
the remaining assertion is
$\sum_{j=1}^{m-1}t^j\geq(m-1)t^m$.
Gronwall's inequality consequently bounds the fundamental solution
$U(0)=I$ on $|z|\leq R<1$ by
\[
 K_R=\exp\!\left[B\left(-\frac{\log(1-R)}m
                    +\frac{m-1}{m}R\right)\right].
\]
Cauchy's estimate bounds the tail after degree $n$ on $|z|\leq r<R$
by
\[
 K_R\frac{(r/R)^{n+1}}{1-r/R}.
\tag{1.18.3}
\]
The code uses $R=0.999$.  The coefficients are computed directly
from (1.18.2):
\[
 nU_n=-(n-m)U_{n-m}
       +\sum_{j=0}^{\min(m-1,n-1)}U_{n-1-j}B_j,
\]
with the first term omitted when $n<m$.
There is no derivative approximation in this recursion.
The regular limit frame is determined by certified continuation
from $\lambda=1$ to a positive root-chart point and inversion of
this convergent series.  At the infinity center it yields
\[
 h_b(\infty)\in
 [-1.094540250311402780116268263054945704
       \ \mathbin{\pm}\ 3\cdot10^{-37}].
\tag{1.18.4}
\]
The corresponding enclosure at zero has center
\[
-0.9317933763565235054307951057960111149908282005165194125
\]
and radius $7.91\cdot10^{-56}$.

We explain two numerical details which are needed to make these
bounds rigorous and efficient.  At a root center, let $H$ denote
the actual limiting period frame, and form the real matrix whose
five columns are
\[
 \operatorname{Im}H_{\mathrm{low},1},\quad
 \operatorname{Re}H_{\mathrm{low},1},\quad
 \operatorname{Im}H_{\mathrm{low},2},\quad
 \operatorname{Re}H_{\mathrm{low},2},\quad
 \operatorname{Im}H_{\mathrm{low},3}.
\]
Its certified inverse gives a constant real change of the lower five
flat coordinates.  Subtracting the corresponding real linear
combination of those coordinates from the top coordinate also
makes the first two top entries zero at the center.  This is an
explicit real invertible change fixing $e$.  It does not change
$h_b$: on the real intersection vector the extra top term is real,
so its imaginary part is zero.  This change greatly improves the
four real equations used to evaluate (1.15.3).

For a polynomial $p(z)=\sum c_jz^j$ with ball coefficients, evaluation
on a rectangle centered at $z_c$, of enclosing radius $\rho$, uses
its midpoint-coefficient polynomial at the point $z_c$.  Coefficient
uncertainty is bounded by
$\sum_j\operatorname{rad}(c_j)|z_c|^j$, and variation over the
rectangle by
$\rho\sum_j j|c_j|r^{j-1}$, where $r$ bounds the modulus on it.
The tail (1.18.3) is then added.  The point evaluation uses 800-digit
ball arithmetic; its rounding radius is retained.  These explicit
bounds avoid the excessive interval enlargement of long rectangular
Horner recursions.  No error radius is discarded.

On the square $|\operatorname{Re}z_\infty|,
|\operatorname{Im}z_\infty|\leq0.1$, differentiating the exact linear
system for the real intersection vector gives the following
certified bounds for the real Hessian of $h_b$:
\[
\begin{split}
 0.2898644&<h_{xx},h_{yy}<0.3946911,\\
 |h_{xy}|&<0.0524256,\qquad
 \det\operatorname{Hess}(h_b)>0.0757786.
\end{split}
\tag{1.18.5}
\]
The derivatives of the period frame are evaluated by
$Y'=YG$ and $Y''=Y(G^2+G')$; the derivatives of the intersection
vector follow by differentiating its invertible four-by-four real
system.  Finite rotation invariance forces $dh_b(0)=0$ in this
root chart.  Positive definiteness on the convex square therefore
proves a strict minimum at its center.

Outside that square, a recursive dyadic rectangle cover proves
$h_b>h_b(\infty)$ throughout $|z_\infty|\leq0.98$ in its sector.
It has 900 accepted rectangles, maximum subdivision depth eight.
The zero-root disc $|z_0|\leq0.98$ is certified by 463 rectangles,
also of maximum depth eight.  A rectangle is discarded only when
its entire intersection with the required sector and disc is empty,
or when it lies in the already proved local square.  Every remaining
rectangle is either certified by the strict interval inequality or
subdivided.  The complete records show that the subdivision terminates;
there is no unresolved rectangle.

It remains to cover a cusp neighborhood and the collar between these
two root discs.  Write
\[
 q=(\lambda-\lambda_*)/\lambda_*.
\]
In this coordinate the connection is $R/q+A(q)$, with $R^2=0$
and $A$ analytic on $|q|<1$.  A regular-singular factorization is
\[
 Y(q)=H\exp((\log q)R)P(q),\qquad P(0)=I.
\]
The recursion for its analytic factor is
\[
 (n-\operatorname{ad}_R)P_n
   =\sum_{j=0}^{n-1}P_{n-1-j}A_j,
 \qquad \operatorname{ad}_R(Q)=QR-RQ.
\]
Because $R^2=0$, the cube of this adjoint operator vanishes and its
inverse is the finite expression
$1/n+\operatorname{ad}_R/n^2+\operatorname{ad}_R^2/n^3$.
In an exact dyadic diagonal norm the code proves
$\|R\|<0.841$ and uses the sharper recorded bound
\[
 C=1+2\|R\|+4\|R\|^2<5.508.
\]
On the circle of radius $0.4$, the absolute-coefficient majorant of
$A$ is less than $4.895$.  The coefficient recursion is dominated
by the scalar equation $p'=Ca_{\rm abs}p$; hence
$\|P\|\leq\exp(C\cdot4.895\cdot0.4)$ there.  Cauchy's estimate
now gives the same explicit geometric tail bound as (1.18.3), with
this exponential in place of $K_R$ and $R=0.4$.

The constant $H$ is obtained by certified continuation to $q=0.01$
along a specified polygon in the upper half-plane.  The matrix
$2\pi iHRH^{-1}$ is integral, since it is the logarithm of the
corresponding marked local monodromy.  Its ball enclosure contains
exactly one integer matrix; that matrix and an integral Jordan
marking $J$ are exported.  They satisfy
\[
 J^{-1}(2\pi iHRH^{-1})J=\begin{pmatrix}0&I_3\\0&0\end{pmatrix}.
\]
At the fixed parameter $c_*=-h_b(\infty)$, the last three columns of
$J^{-1}HP(q)$ specifying $F_{c_*}$ have blocks $T(q),D(q)$.
The lower block is certified invertible on the stated cusp disc.
The period matrix is therefore
\[
 Z(q)=\frac{\log q}{2\pi i}I_3+C(q),\qquad C(q)=T(q)D(q)^{-1}.
\]
Put $B(q)=\operatorname{Im}C(q)$ and $y=-\log|q|/(2\pi)$.
Then
\[
 \det\operatorname{Im}Z
 =y^3+\operatorname{tr}(B)y^2
   +\Bigl(\sum_{i<j}(B_{ii}B_{jj}-B_{ij}B_{ji})\Bigr)y+\det B.
\tag{1.18.6}
\]
On the whole disc $|q|\leq0.01$, interval substitution in (1.18.6),
followed by $y=-\log(0.01)/(2\pi)+v$, gives four strictly positive
coefficients in $v$; their certified lower bounds, in increasing
order, are $0.8111$, $3.6097$, $3.5346$, and $1$.
The sharp endpoints are in \texttt{cusp\_cubic\_bounds.json}.
Since $v\geq0$, the determinant is positive
throughout the punctured disc.  Near the cusp $h_b+c_*>0$ by P1.8;
continuity and transversality then give this same strict inequality
throughout that disc.  A separate finite cover of
$0.01\leq|q|\leq0.25$, using the same analytic factor and (1.18.6),
contains 7033 certified rectangles, of maximum depth seven.  It
extends the inequality to the whole punctured disc $|q|\leq0.25$.

The only remaining region, expressed in the infinity-root chart, is
\[
 0.98\leq |z_\infty|\leq0.98^{-4/5},\qquad
 0\leq\arg z_\infty\leq2\pi/5,
 \qquad |q|\geq0.25,
\]
where $q=-z_\infty^{-5}-1$.
Each rectangle in this collar is reached from radius $0.95$ by
Taylor continuation of (1.18.2).  The local analytic radius is
$0.449$ times the distance to the nearest root of $1+z^5$.
The numerator is bounded in the same diagonal norm and the
denominator by the product of the five distances to those roots.
Gronwall and Cauchy give an explicit tail for every step, exactly
as in P1.16.  There are 427 steps and 117 accepted rectangles,
of maximum depth eight; each proves $h_b>h_b(\infty)$.

These regions cover every point of the completed base except the
cusp itself, where $h_b$ tends to positive infinity.  Every inequality
away from the infinity center is strict.  This proves both the
global minimum and its uniqueness, and P1.8 gives the exact parameter
domain.

Finally, \texttt{normalization\_shift.py} evaluates
$a_{\rm split}=96Y_{03}(\lambda_*/2)/(625\pi i)$ by the same
certified continuation, with the polygon stated above.
The real and imaginary radii are respectively less than
$4.31\cdot10^{-56}$ and $3.15\cdot10^{-56}$.
Subtracting its imaginary part from (1.18.4) gives the stated
point-normalized minimum and interval.  This is the explicit
splitting morphism already established in P1.15, now numerically
certified.
\end{proof}

\begin{statement}{P1.19}{Proved}
There is an independent computation from the full global cellular
complex.  Its chain ranks in degrees zero through eight are
\[
 (4,26,87,162,200,156,79,23,3),
\]
and its differentials, after the explicit integral chain equivalences
below, are the matrices of case one in
\path{code/part1/global_ext_cases.json}.
Their integral Smith calculation gives
\[
 H_k(M;\mathbb Z)=
 \begin{cases}
 \mathbb Z,&k=0,8,\\
 \mathbb Z^4,&k=4,\\
 0,&\text{otherwise}.
 \end{cases}
\]
This calculation uses local Poincar\'e--Lefschetz duality of the
smooth cusp cap to determine its secondary attaching class.  It does
not use global Poincar\'e duality or the homology conclusion of P1.12.
In particular, the map obtained by setting the whole cusp suspension
chain term to zero is not the actual relative cusp model; its precise
missing term is accounted for here.
\end{statement}
\begin{proof}
We retain the actual complement complex and finite-cap maps of P1.12.
The issue left by that argument was the full cusp boundary map.
We first determine its homology maps, then its secondary chain term,
and finally give the explicit integral reduction used by the program.

Write the cusp lattice as $A\oplus B$, with ordered generators
$a_1,a_2,a_3,b_1,b_2,b_3$ and $Nb_i=a_i$.
The boundary is the mapping torus of
\[
 (b,a)\longmapsto(b,a+b)
\]
on $T_B^3\times T_A^3$.
For the pure toric model, its retraction to the central moment model
has the following explicit description.  On the dual cell of a Kuhn
simplex $\sigma$, choose any vertex $v$ of $\sigma$ and put
\[
 (b,a,\theta)\longmapsto
   [b,a+\theta(b-v)],\qquad0\leq\theta\leq1.
\tag{1.19.1}
\]
The brackets denote the angular quotient by $L_\sigma$.
Changing $v$ changes the displayed angular coordinate by a real
multiple of a vector in $L_\sigma$, so the map is unchanged.
Replacing $b$ by $b+n$ also replaces $v$ by $v+n$; hence the formula
is periodic.  At $\theta=1$ the integer translation $v$ disappears
on the angular torus, and the map agrees with the cusp monodromy
identification.  These observations prove that (1.19.1) is well-defined
on the boundary mapping torus.  In a monomial chart of P1.9 it is
exactly the operation of forgetting the arguments of the vanishing
coordinates, with the argument of $q$ recorded by $\theta$.
Thus it is the actual boundary retraction.  The deformation removing
the bounded period units in P1.11 transports this map by a homotopy
and preserves its homology class.

We prove that the suspension classes have zero image in cusp homology.
Let $D$ be the derivation of $\bigwedge^\bullet(A\oplus B)$ extending
$Nb_i=a_i$, $Na_i=0$.  Then $\bigwedge M_0=\exp D$ and
$\ker(\bigwedge M_0-I)=\ker D$ over $\mathbb Q$.
An integral vector in this kernel has the same property.
The summand with $p$ base generators and $q$ angular generators is
sent by $D$ to the summand with $(p-1,q+1)$; consequently each
homogeneous summand of an invariant vector is separately killed.
Let $E$ replace $a_i$ by $b_i$. In the orthonormal exterior basis
$E=D^t$ and $[D,E]=(q-p)I$ on that summand.  For $Dv=0$,
\[
 (q-p)\|v\|^2
 =\langle[D,E]v,v\rangle=\|Ev\|^2\geq0.
\]
Thus an invariant $k$-cycle has base degree at most $\lfloor k/2\rfloor$.

This assertion about cycles can be realized geometrically, rather
than only in a rational associated graded group.  Give $T_B^3$ its
minimal product CW structure and put
$F_p=T_A^3\times(T_B^3)^{(p)}$.
Its cellular differential is zero, so
$H_*(F_p;\mathbb Z)\to H_*(T^6;\mathbb Z)$ is injective.
The shear preserves $F_p$ and admits cellular approximation through
maps preserving this filtration.  The corresponding mapping-cylinder
cells give the integral complex
$\operatorname{Cone}(\bigwedge M_0-I)$.
An invariant cellular vector of base degree at most $p$ therefore
has a suspension lift whose entire projection to $T_B^3$ has dimension
at most $p$.  Equivalently, its closing boundary already bounds in
$F_p$, by the just-proved injectivity.  This construction fixes the
suspension representatives used in the Wang complex and is compatible
with the based product-cell comparison in P1.12.

Formula (1.19.1) preserves $b$, so the image of such a lift also has
base projection of dimension at most $p$.  It represents filtration
at most $p$ in the dual-cell filtration of P1.11.  Here is a direct
cellular justification of this last step.  In every dual cell of
dimension greater than $p$, choose an interior point outside that
projection.  Radial retraction from this point pushes the projected
cycle to the cell boundary.  On the central moment model it lifts
by keeping the angular coordinate fixed; on the boundary the angular
quotient is the canonical further quotient.  The lift is continuous
and fixes the boundary.  Performing these retractions in descending
cell dimension gives a homologous cycle on the dual $p$-skeleton.
There are finitely many cells modulo periods, so this operation
terminates.

The integral $E_2=E_\infty$ table of P1.11 has nonzero entries only
at angular degree $q\leq d$, where $d$ is the base-cell degree.
For a suspension of a $k$-cycle, the total degree is $k+1$, whereas
$d\leq p\leq\lfloor k/2\rfloor$. Its angular degree would therefore
be $k+1-d>d$, where that table is zero.  The suspension image
vanishes.  On the fibre summand the map is precisely the marked
collapse $\kappa_k$ already exported in
\texttt{cusp\_marked\_collapse.json}.  Thus the homology maps agree
with those of the chain map $f_0=(\kappa,0)$, in every degree.

Agreement on homology does not identify these two chain maps.
We now determine the remaining difference.  Integral Smith reduction
splits a finite free complex into free homology generators, contractible
unit pairs, and two-term pairs with differential equal to a torsion
invariant factor.  The cusp homology is free.  The boundary homology
has just one torsion group, $\mathbb Z/2$ in degree three.
Consequently, after maps on homology have been fixed, the difference
of two maps from the boundary complex to the cusp complex is
represented by one vector
\[
 \varepsilon\in H_4(C;\mathbb Z)/2H_4(C;\mathbb Z)
                 \cong(\mathbb Z/2)^{10}.
\tag{1.19.2}
\]
This follows directly by applying a map to the upper generator of
the two-term pair $\mathbb Z\xrightarrow{2}\mathbb Z$; changing a
chain homotopy changes its image by twice an integral vector.
There are no other secondary terms.  Thus no rational formality or
unproved equivariant formality is used at this point.

In the marked cusp degree-four basis, the primitive functional
\[
 \psi=(0,0,0,1,-1,1,1,-1,-1,1)
\tag{1.19.3}
\]
annihilates $\kappa_4$ and identifies its cokernel with $\mathbb Z$.
The exact multiplication $\psi\kappa_4=0$ and the primitive Smith
factors are exported.  Setting $\varepsilon=0$ gives relative
cusp homology $H_4(C,\partial)=\mathbb Z^{10}\oplus\mathbb Z/2$,
as computed in \texttt{cusp\_relative\_candidate.json}.
More intrinsically, the relative exact sequence in this degree is
an extension of $\mathbb Z^9\oplus\mathbb Z/2$ by the cokernel
$\mathbb Z$ detected by $\psi$.  Its torsion extension class is
exactly $\psi\varepsilon$ modulo two.

The actual cusp neighborhood is a compact oriented smooth
8-manifold with boundary: take the inverse image of a sufficiently
small closed disc, whose boundary circle avoids the critical value.
It retracts onto $C$.  Poincar\'e--Lefschetz duality and the already
computed free cusp homology give
\[
 H_4(\text{cusp cap},\partial;\mathbb Z)
       \cong H^4(C;\mathbb Z)=\mathbb Z^{10}.
\]
The preceding extension must therefore be nonsplit on its torsion
summand.  This determines the required secondary class for the global
calculation:
\[
                 \psi\varepsilon=1\pmod2.
\tag{1.19.4}
\]
This also proves the stated failure of the wholly zero suspension
chain map as a relative cusp model.

We next give the integral equivalence which reduces every class
satisfying (1.19.4) to one displayed matrix.  Write the degree-four
homology gluing map in the plus-sign convention as
\[
 j_4:H_4(\partial)\longrightarrow H_4(Y)\oplus H_4(C).
\]
Multiplication by $-1$ on the cusp chain summand is the chain
isomorphism from the minus-sign convention in the exported whole
diagram to this convention.  The matrix $j_4$ is $22$ by $18$,
has rank eighteen, and all its nonzero Smith factors are one.
Its cokernel is free of rank four. In its exported cokernel basis,
the map from $H_4(C)$ is
\[
 \begin{pmatrix}0\\0\\0\\\psi\end{pmatrix}.
\tag{1.19.5}
\]
This is a direct integer matrix calculation, independent of the
homology of $M$.

Let $e_4$ denote the fourth coordinate vector of $\mathbb Z^{10}$,
so $\psi e_4=1$.  The nine columns
$e_i-\psi_i e_4$, $i\ne4$, form an integral basis $K$ of $\ker\psi$;
$[e_4,K]$ is unimodular.  The file
\texttt{cusp\_ext\_reduction.json} gives integral lifts $W$ satisfying
\[
 j_4W=\begin{pmatrix}0\\K\end{pmatrix}.
\]
It also gives actual boundary cycles $X$ representing these lifts,
and integral complement chains $T$ with
\[
 d_YT=j_YX,\qquad f_0X=K.
\tag{1.19.6}
\]
Every entry is checked by multiplication in the program.

For clarity, the chain automorphisms implementing the reduction
are equally explicit.  Let $r$ be the exported integral degree-four
cocycle representing the torsion generator in $H^4(\partial;\mathbb Z)$.
In the ordered 35-dimensional boundary chain group it is
$r=-e_{17}^{t}-e_{18}^{t}$.  It satisfies
\[
 r d_{\partial,5}=0,\qquad rX=0,
 \qquad 2r=t\,d_{\partial,4},
\tag{1.19.7}
\]
where the integral row $t$ is exported as well.
For a column $v\in\mathbb Z^9$, the endomorphism
$\eta=(Xv)r$ of degree-four boundary chains, zero in all other
degrees, is a chain map and has $\eta^2=0$.
Thus $I+\eta$ is an integral chain automorphism with inverse
$I-\eta$. Precomposing the attaching map by it changes the
secondary vector by $Kv$, up to the explicitly given homotopy
$(Tv)r$ in the complement, by (1.19.6).
The last equation in (1.19.7) removes an even multiple of $e_4$
by the homotopy $e_4t$.  Decomposing $\varepsilon$ in the
unimodular basis $[e_4,K]$, equation (1.19.4) therefore reduces
the actual global attaching map to the canonical correction
\[
 f_4=f_{0,4}+e_4r,
 \qquad f_k=f_{0,k}\quad(k\ne4).
\tag{1.19.8}
\]
These formulas prove the chain equivalence; they do not merely
compare the ranks of the two complexes.

The program \texttt{global\_ext\_cases.py} inserts (1.19.8) into
the complete main-piece/three-cap mapping cone, retaining every
finite-cap differential and every boundary map of P1.12.
It checks $d^2=0$ over $\mathbb Z$ and exports every differential.
Its odd case is the actual global model by the preceding argument.
The even case is also computed as a diagnostic, and happens to give
the same global homology; it is not used to identify the local cap.
The independent program \texttt{global\_cell\_certificate.py}
exports the Smith left and right matrices, kernel bases, incoming
matrices in those bases, and their Smith left and right matrices
for the odd case.  The differential ranks in degrees one through
eight are
\[
                 (3,23,64,98,98,58,21,2).
\]
All nonzero incoming Smith factors in the cycle lattices are one.
Subtracting incoming and outgoing ranks from the displayed chain
ranks gives $(1,0,0,0,4,0,0,0,1)$, with no torsion.
This completes the full cellular comparison without using the
global duality conclusion of P1.12.
\end{proof}
